\documentclass[
10pt, % Main document font size
a4paper, % Paper type, use 'letterpaper' for US Letter paper
oneside, % One page layout (no page indentation)
%twoside, % Two page layout (page indentation for binding and different headers)
headinclude,footinclude, % Extra spacing for the header and footer
BCOR5mm, % Binding correction
]{scrartcl}
%%%%%%%%%%%%%%%%%%%%%%%%%%%%%%%%%%%%%%%%%
% Arsclassica Article
% Structure Specification File
%
% This file has been downloaded from:
% http://www.LaTeXTemplates.com
%
% Original author:
% Lorenzo Pantieri (http://www.lorenzopantieri.net) with extensive modifications by:
% Vel (vel@latextemplates.com)
%
% License:
% CC BY-NC-SA 3.0 (http://creativecommons.org/licenses/by-nc-sa/3.0/)
%
%%%%%%%%%%%%%%%%%%%%%%%%%%%%%%%%%%%%%%%%%

%----------------------------------------------------------------------------------------
%	REQUIRED PACKAGES
%----------------------------------------------------------------------------------------

\usepackage[
nochapters, % Turn off chapters since this is an article        
beramono, % Use the Bera Mono font for monospaced text (\texttt)
eulermath,% Use the Euler font for mathematics
pdfspacing, % Makes use of pdftex’ letter spacing capabilities via the microtype package
dottedtoc % Dotted lines leading to the page numbers in the table of contents
]{classicthesis} % The layout is based on the Classic Thesis style

\usepackage{arsclassica} % Modifies the Classic Thesis package
\usepackage{ulem}
\usepackage{graphicx}
\usepackage{float}
\usepackage{algorithm2e}
\usepackage[export]{adjustbox}
\usepackage[utf8]{inputenc}  
\usepackage[english,frenchb]{babel} 
\usepackage[T1]{fontenc} % Required for accented characters
\usepackage{multicol}
\usepackage{graphicx} % Required for including images
\graphicspath{{Figures/}} % Set the default folder for images

\usepackage{enumitem} % Required for manipulating the whitespace between and within lists

\usepackage{lipsum} % Used for inserting dummy 'Lorem ipsum' text into the template


\expandafter\def\csname ver@subfig.sty\endcsname{}


\usepackage{svg}
\usepackage{hyperref}

\usepackage{amsmath,amssymb,amsthm} % For including math equations, theorems, symbols, etc

\usepackage{varioref} % More descriptive referencing

%----------------------------------------------------------------------------------------
%	THEOREM STYLES
%---------------------------------------------------------------------------------------

\theoremstyle{definition} % Define theorem styles here based on the definition style (used for definitions and examples)
\newtheorem{definition}{Definition}

\newtheorem*{remark}{Remark}
\theoremstyle{plain} % Define theorem styles here based on the plain style (used for theorems, lemmas, propositions)
\newtheorem{theorem}{Theorème}
\newtheorem{lemma}[theorem]{Lemme}
\newtheorem{proposition}[theorem]{Proposition}

\theoremstyle{remark} % Define theorem styles here based on the remark style (used for remarks and notes)

\DeclareMathOperator*{\argmin}{arg\,min}
\DeclareMathOperator*{\argmax}{arg\,max}

%----------------------------------------------------------------------------------------
%	HYPERLINKS
%---------------------------------------------------------------------------------------

\hypersetup{
%draft, % Uncomment to remove all links (useful for printing in black and white)
colorlinks=true, breaklinks=true, bookmarks=true,bookmarksnumbered,
urlcolor=webbrown, linkcolor=RoyalBlue, citecolor=webgreen, % Link colors
pdftitle={}, % PDF title
pdfauthor={\textcopyright}, % PDF Author
pdfsubject={}, % PDF Subject
pdfkeywords={}, % PDF Keywords
pdfcreator={pdfLaTeX}, % PDF Creator
pdfproducer={LaTeX with hyperref and ClassicThesis} % PDF producer
}

\hyphenation{Fortran hy-phen-ation} 

\title{\normalfont{FTML Exercices 4 solutions}\\
} 
\author{\spacedlowsmallcaps{}
} 
\date{\today}
\begin{document}
\renewcommand{\sectionmark}[1]{\markright{\spacedlowsmallcaps{#1}}}
\lehead{\mbox{\llap{\small\thepage\kern1em\color{halfgray} \vline}\color{halfgray}\hspace{0.5em}\rightmark\hfil}} 
\pagestyle{scrheadings}
\maketitle 
\setcounter{tocdepth}{3} 

\tableofcontents

\section{\large\color{Blue}Forte convexité de la loss Ridge regression}

\textbf{1)} On veut montrer que pour $\alpha\in [0,1]$, et $x, y\in \mathbb{R}$:

\begin{equation}
    (\alpha x + (1-\alpha)y)^2\leq \alpha x^2 + (1-\alpha)y^2 - \alpha(1-\alpha)(x-y)^2
    \label{eq:}
\end{equation}

C'est-à-dire:
\begin{equation}
    \alpha^2 x^2 + (1-\alpha)^2y^2+2\alpha(1-\alpha)xy\leq \alpha x^2 + (1-\alpha)y^2 - \alpha(1-\alpha)(x-y)^2
    \label{eq:}
\end{equation}

En mettant tous les termes du même côté, cela revient à:

\begin{equation}
    \alpha x^2-\alpha^2 x^2 + (1-\alpha)y^2 -(1-\alpha)^2y^2- \alpha(1-\alpha)(x-y)^2-2\alpha(1-\alpha)xy \geq 0
    \label{eq:}
\end{equation}

On remarque que
$\alpha-\alpha^2=\alpha(1-\alpha)$
et que 
$(1-\alpha)-(1-\alpha)^2=\alpha(1-\alpha)$

Cela permet de tout factoriser par $\alpha(1-\alpha)$. Il nous reste donc à vérifier que:

\begin{equation}
     x^2 + y^2- (x-y)^2-2xy \geq 0
    \label{eq:}
\end{equation}

Or ceci est effectivement vrai car:

     \begin{align}
         x^2 + y^2- (x-y)^2-2xy
         &= x^2 + y^2- (x-y)^2-2xy \\
         &= x^2 + y^2- x^2-y^2+2xy-2xy \\
         &= 0\\
     \end{align}

On est donc dans un exemple où l'inégalité est saturée.
\\

\textbf{2)} On remarque que pour $x\mapsto||x||^2$,  avec $x\in \mathbb{R}^d$, les premières étapes du calcul sont identiques. En effet, on peut utiliser le fait que pour tout vecteur $z$, $||z||^2=\langle z, z \rangle$ et remplacer:
\begin{itemize}
    \item les $x^2$ par des $||x||^2$
    \item les $y^2$ par des $||y||^2$
    \item les $xy$ par des $\langle x, y \rangle$ (produit scalaire)
\end{itemize}

Il nous reste donc à vérifier que:

\begin{equation}
    ||x||^2 + ||y||^2- ||x-y||^2-2\langle x, y\rangle \geq 0
    \label{eq:}
\end{equation}

Or c'est aussi égal à $0$ en développant comme précédemment.
\\

\textbf{3)} Avec la question 2), on sait que l'application $\theta \mapsto \lambda ||\theta||^2$ est $2\lambda$ fortement convexe, dès que $\lambda>0$. Avec les notation suivantes:
\begin{itemize}
    \item $u:\theta \mapsto \frac{1}{n}||X\theta -y||^2$, qui est convexe.
    \item $v:\theta \mapsto \lambda||\theta||^2$, qui est donc $2\lambda$ fortement convexe.
\end{itemize}

on note que la loss Ridge est égale à $u(\theta)+v(\theta)$. Elle est donc également $2\lambda$ fortement convexe. 

En effet:

\begin{align}
    (u+v)(\alpha\theta_1+(1-\alpha)\theta_2)&=
    u(\alpha\theta_1+(1-\alpha)\theta_2)+v(\alpha\theta_1+(1-\alpha)\theta_2)\\
    &\leq \alpha u(\theta_1)+(1-\alpha)u(\theta_2)+\alpha v(\theta_1)+(1-\alpha)v(\theta_2)-\alpha(1-\alpha)||\theta_1-\theta_2||^2\\
    &\leq \alpha (u+v)(\theta_1)+(1-\alpha)(u+v)(\theta_2) -\alpha(1-\alpha)||\theta_1-\theta_2||^2\\
 \end{align} 

\textbf{4)} Soit $\mu = 2\lambda_{min}$. Montrons que:

\begin{equation}
    (\alpha x+(1-\alpha)y)^TQ(\alpha x+(1-\alpha)y)\leq \alpha x^TQx+(1-\alpha)y^TQy- \frac{\mu}{2} \alpha(1-\alpha)||x-y||^2
\end{equation}

c'est-à-dire:

\begin{equation}
    (\alpha x+(1-\alpha)y)^TQ(\alpha x+(1-\alpha)y) -\alpha x^TQx-(1-\alpha)y^TQy\leq - \frac{\mu}{2} \alpha(1-\alpha)||x-y||^2
\end{equation}

We compute the left-hand side:

\begin{equation}
    \begin{aligned}
        \label{eq:}
	&(\alpha x+(1-\alpha)y)^TQ(\alpha x+(1-\alpha)y) -\alpha x^TQx+(1-\alpha)y^TQy \\
	&= \alpha^2 x^TQx + (1-\alpha)^2 y^TQy + \alpha(1-\alpha)(x^TQy+y^TQx)-\alpha x^TQx-(1-\alpha)y^TQy\\
	&= \alpha(\alpha-1) x^TQx + (1-\alpha)((1-\alpha)-1) y^TQy + \alpha(1-\alpha)(x^TQy+y^TQx)\\
	&= \alpha(\alpha-1) x^TQx + \alpha(\alpha-1) y^TQy + \alpha(1-\alpha)(x^TQy+y^TQx)\\
	&= \alpha(1-\alpha) \Big(-x^TQx -  y^TQy + x^TQy+y^TQx\Big)\\
	&= -\alpha(1-\alpha) \Big((x-y)^TQ(x-y)\Big)\\
	&\leq -\lambda_{min}\alpha(1-\alpha)||x-y||^2
    \end{aligned}
\end{equation}

Pour arriver à la dernière ligne, on a utilisé le résultat suivant: pour tout vecteur $x\in \mathbb{R}^d$, $x^TQx=\langle Qx, x\rangle\geq \lambda_{min} ||x||^2$. On peut prouver ce résultat en décomposant $x$ dans une base orthonormée de vecteurs propres de $Q$.

\end{document}
